% Blueprint of the Lean development (lean/ReflectZphi). Statement numbers refer to the paper
% "Rank 3 reflective hyperbolic lattices over real quadratic fields". The table
% ../reflect-zphi-paper/lean-map.md is the source of the correspondence.

\chapter*{About this blueprint}

This blueprint lists the statements of the paper that are formalised in Lean~4 with Mathlib
(the library \texttt{ReflectZphi} in the folder \texttt{lean/} of the repository), and shows
how they depend on each other. A statement marked as formalised is a Lean definition or
theorem; its name is given and links to the source. The results imported without proof
(Allcock's thin-part theorem, Vinberg's chamber theorem) are Lean definitions of
propositions, taken as explicit hypotheses by the theorems that use them; they are not axioms.

The main theorem is Theorem~\ref{thm:all-reflective}: the $144$ classes over $\ZZ[\varphi]$
are reflective. It is proved by kernel-checked certificates, one per class, each an
application of the certificate checker (Theorems~\ref{thm:reflective-of-check}
and~\ref{thm:reflective-of-checkM}). Chapter~\ref{chap:not-formalised} lists what is proved
only in the paper.

\chapter{The rings}\label{chap:rings}

The rings are $\ZZ[\omega]$ with $\omega^2=a+b\omega$, as \texttt{QuadraticAlgebra ℤ a b}:
$\ZZ[\varphi]$ is $(a,b)=(1,1)$, $\ZZ[\sqrt2]$ is $(2,0)$, and $\ZZ[(1+\sqrt{13})/2]$ is $(3,1)$.

\begin{definition}[real embeddings]\label{def:emb}
\lean{ReflectZphi.emb, ReflectZphi.emb_apply}\leanok
For $b^2+4a>0$ the two real embeddings $\iota_{1,2}$ send $\omega$ to $(b\pm\sqrt{b^2+4a})/2$.
\end{definition}

\begin{lemma}[exact sign test]\label{lem:sign}
\lean{ReflectZphi.posCode, ReflectZphi.pos_iff_posCode, ReflectZphi.emb_pos_iff}\leanok
\uses{def:emb}
The sign of $\iota_j(x+y\omega)$ is decided by integer comparisons on $x$, $y$, $a$, $b$.
\end{lemma}
\begin{proof}\leanok
Compare squares; paper, Section~9.
\end{proof}

\begin{theorem}[Euclidean, hence PID; paper, Prop.~9.1(1)]\label{thm:euclid}
\lean{ReflectZphi.euclideanDomain, ReflectZphi.isPrincipalIdealRing}\leanok
If $|a|+|b|\le2$ and $b^2+4a$ is not a square, then $\ZZ[\omega]$ is norm-Euclidean, hence
a principal ideal domain.
\end{theorem}
\begin{proof}\leanok
Division with rounding (paper, Lemma~2.3).
\end{proof}

\begin{theorem}[$\ZZ\lbrack\varphi\rbrack$; paper, Prop.~9.1]\label{thm:zphi}
\lean{ReflectZphi.zphi_isPID, ReflectZphi.iota1_phi, ReflectZphi.iota2_phi, ReflectZphi.norm_phi, ReflectZphi.isUnit_phi, ReflectZphi.five_eq, ReflectZphi.norm_piZ, ReflectZphi.norm_mod_five, ReflectZphi.prime_two, ReflectZphi.prime_three}\leanok
\uses{thm:euclid, def:emb}
$\ZZ[\varphi]$ is a PID; $\iota_1(\varphi)$ is the golden ratio and $\iota_2(\varphi)$ its
conjugate; $N(\varphi)=-1$, so $\varphi$ is a unit; $5=\varphi^2\pi^2$ with $N(\pi)=5$; norms
are $0,\pm1$ modulo $5$; and $2$ and $3$ are prime.
\end{theorem}
\begin{proof}\leanok
Theorem~\ref{thm:euclid} and direct computation.
\end{proof}

\begin{theorem}[$\ZZ\lbrack\sqrt2\rbrack$; paper, Section~8]\label{thm:zsqrt2}
\lean{ReflectZphi.zsqrt2_isPID, ReflectZphi.norm_epsSqrt2, ReflectZphi.iota1S_sqrt2}\leanok
\uses{thm:euclid}
$\ZZ[\sqrt2]$ is a PID, $N(1+\sqrt2)=-1$, and $\iota_1(\sqrt2)=\sqrt2$.
\end{theorem}
\begin{proof}\leanok
Theorem~\ref{thm:euclid}.
\end{proof}

\begin{theorem}[$\ZZ\lbrack\omega\rbrack$, $\omega=(1+\sqrt{13})/2$; paper, Prop.~13.1]\label{thm:zsqrt13}
\lean{ReflectZphi.boxBound_three_one, ReflectZphi.isPrincipalIdealRing_of_boxBound, ReflectZphi.zsqrt13_isPID, ReflectZphi.norm_eps13, ReflectZphi.isUnit_eps13, ReflectZphi.three_eq, ReflectZphi.thirteen_eq, ReflectZphi.norm_ne_two, ReflectZphi.prime_two_13}\leanok
Division with remainder holds for $(a,b)=(3,1)$, so the ring is a PID; $N(1+\omega)=-1$;
$3=-\omega(1-\omega)$ and $13=(2\omega-1)^2$; and $2$ is prime.
\end{theorem}
\begin{proof}\leanok
A box bound for the rounding error, then direct computation.
\end{proof}

\chapter{Lattices and the imported results}\label{chap:lattices}

\begin{definition}[admissible lattices, roots, reflections; paper, Section~2]\label{def:lattice}
\lean{ReflectZphi.Lattice.Admissible, ReflectZphi.Lattice.Primitive, ReflectZphi.Lattice.IsRoot, ReflectZphi.Lattice.IsReflectionOf}\leanok
\uses{def:emb}
An admissible lattice is $\Oo^3$ with a symmetric Gram matrix $G$ of signature $(2,1)$ at
$\iota_1$ and positive definite at $\iota_2$. A root is a primitive $r$ with $r^2$ totally
positive and $2(r\cdot x)/r^2\in\Oo$ for all $x$; its reflection is
$x\mapsto x-2(r\cdot x)/r^2\,r$.
\end{definition}

\begin{definition}[reflective]\label{def:reflective}
\lean{ReflectZphi.Lattice.isometryGroup, ReflectZphi.Lattice.reflectionGroup, ReflectZphi.Lattice.Reflective}\leanok
\uses{def:lattice}
$O(L)$ is the group of isometries of $L$, and $W(L)$ is generated by the reflections in the
roots. $L$ is reflective if $W(L)$ has finite index in $O(L)$.
\end{definition}

\begin{definition}[Vinberg's chamber theorem, imported; paper, App.~A.1(2)]\label{def:vinberg}
\lean{ReflectZphi.VinbergChamber}\leanok
\uses{def:reflective}
The proposition: for an admissible lattice over $\ZZ[\omega]$ with non-square discriminant,
$W(L)$ acts simply transitively on its chambers in either sheet, and $L$ is reflective if
and only if some chamber is compact \cite{Vinberg}. It is a hypothesis of every
reflectivity theorem, not an axiom.
\end{definition}

\begin{definition}[Allcock's thin-part theorem, imported; paper, Thm.~5.1, App.~A.1(1)]\label{def:allcock}
\lean{ReflectZphi.AllcockThinPart}\leanok
\uses{def:lattice}
The proposition of Allcock's thin-part theorem for a real symmetric form of signature
$(2,1)$. No certificate uses it; the enumeration that does is not formalised
(Chapter~\ref{chap:not-formalised}).
\end{definition}

\chapter{The reflectivity certificates}\label{chap:certificates}

The item numbers (1)--(6) are those of the paper, Appendix~A.2.

\begin{lemma}[Minkowski inequalities, (2)]\label{lem:minkowski}
\lean{ReflectZphi.Minkowski.rev_cs, ReflectZphi.Minkowski.majorant_bound, ReflectZphi.Minkowski.tau_add, ReflectZphi.Minkowski.ball_convex}\leanok
Reverse Cauchy--Schwarz, the bound by a positive majorant, the reverse triangle inequality,
and convexity of hyperbolic balls, for $\operatorname{diag}(1,1,-1)$.
\end{lemma}
\begin{proof}\leanok
Real analysis on $\RR^3$.
\end{proof}

\begin{lemma}[polygons, (2)]\label{lem:polygon}
\lean{ReflectZphi.Minkowski.hull, ReflectZphi.Minkowski.isCompact_polySet, ReflectZphi.Minkowski.nonpos_on_polySet, ReflectZphi.Minkowski.crossing_bound}\leanok
\uses{lem:minkowski}
The polygon is the hull of its vertices and is compact; a vector with one sign at the
vertices has that sign on the polygon; and a mirror crossing the polygon has a root of
bounded height.
\end{lemma}
\begin{proof}\leanok
Convexity and Lemma~\ref{lem:minkowski}.
\end{proof}

\begin{lemma}[no crossing mirror gives a chamber, (2)]\label{lem:chamber}
\lean{ReflectZphi.Minkowski.polySet_eq_chamber, ReflectZphi.Transfer.chamber_of_J}\leanok
\uses{lem:polygon}
If no mirror crosses the interior of the polygon, the polygon is the closure of a connected
component of the sheet minus the mirrors. This transfers from $\operatorname{diag}(1,1,-1)$
coordinates to $G$.
\end{lemma}
\begin{proof}\leanok
Connectedness of the interior and Lemma~\ref{lem:polygon}.
\end{proof}

\begin{theorem}[the bridge to reflectivity, (2)]\label{thm:bridge}
\lean{ReflectZphi.Bridge.isChamber_of_cert, ReflectZphi.Bridge.reflective_of_cert}\leanok
\uses{lem:chamber, def:vinberg, def:reflective}
A polygon certificate and a one-sign certificate give a compact chamber, and so, assuming
Vinberg's chamber theorem, a reflective lattice.
\end{theorem}
\begin{proof}\leanok
Lemma~\ref{lem:chamber} and Definition~\ref{def:vinberg}.
\end{proof}

\begin{lemma}[root norms and the unit window, (3); paper, Lemma~3.2]\label{lem:rootnorms}
\lean{ReflectZphi.RootNorms.root_norm_dvd, ReflectZphi.RootNorms.exists_unit_window, ReflectZphi.RootNorms.exists_root_in_window, ReflectZphi.RootNorms.windowNorms, ReflectZphi.RootNorms.winRow, ReflectZphi.RootNorms.winRows, ReflectZphi.RootNorms.windowNorms_of_chunks}\leanok
\uses{def:lattice, thm:zphi}
The norm of a root divides $2\det G$; some unit multiple of a root has its norm $m$ in the
window $\sigma(m)\le m<\frac pq\sigma(m)$ (for $p/q\ge\varepsilon^2/\sigma(\varepsilon)^2$); and the
norms in the window form an explicit finite list, which may be computed row by row.
\end{lemma}
\begin{proof}\leanok
Paper, Section~3 and Appendix~A.2(3).
\end{proof}

\begin{lemma}[enumeration, (5)]\label{lem:enum}
\lean{ReflectZphi.Enum.enum_sound, ReflectZphi.Enum.enumS_sound, ReflectZphi.Enum.quadCheck_sound, ReflectZphi.QuadSolve.solveQuad, ReflectZphi.QuadSolve.mem_solveQuad, ReflectZphi.PolyEnum.hone_of_enum2, ReflectZphi.PolyEnum.reflective_of_enum2}\leanok
\uses{lem:rootnorms}
A Fincke--Pohst enumeration over $\ZZ^n$, certified by a weighted sum of squares, lists
every integer vector in an ellipsoid; the last $\Oo$-coordinate can be solved exactly.
\end{lemma}
\begin{proof}\leanok
The certificate inequalities, checked by the kernel; the iteration itself is not trusted.
\end{proof}

\begin{lemma}[chains of cells, (4)]\label{lem:cells}
\lean{ReflectZphi.Cells.segment_lat, ReflectZphi.Cells.cell_oneSign, ReflectZphi.Cells.hone_of_cells}\leanok
\uses{lem:polygon, lem:enum}
A mirror crossing a segment of the polygon meets a cell of the subdivision; if every cell has
the one-sign property, so has the polygon.
\end{lemma}
\begin{proof}\leanok
Paper, Appendix~A.2(4).
\end{proof}

\begin{lemma}[fast leaves, split trees, per-norm sublattices, (5)--(6)]\label{lem:split}
\lean{ReflectZphi.FastLeaf.leaf3_sound, ReflectZphi.FastLeaf.leaf4_sound, ReflectZphi.FastLeaf.cellCheck_fast4, ReflectZphi.Split.coverAux, ReflectZphi.Split.enumS_of_cover, ReflectZphi.Split.cellCheck_of, ReflectZphi.Split.coverL, ReflectZphi.Split.enumS_of_coverL, ReflectZphi.Sublat.SubCell, ReflectZphi.Sublat.subCheck, ReflectZphi.Sublat.sub_coords, ReflectZphi.Sublat.sub_oneSign, ReflectZphi.Sublat.cellOK_of, ReflectZphi.Sublat.hone_of_cellsM}\leanok
\uses{lem:cells, lem:enum}
Faster forms of the cell check: a table of products at the leaves, a search tree split into
separately checked subtrees glued by a cover of its top, and searches in the sublattices
$L_m=\{r: m\mid 2(Gr)_i\}$ of fixed root norm.
\end{lemma}
\begin{proof}\leanok
Each reduces to Lemma~\ref{lem:cells}.
\end{proof}

\begin{theorem}[the checker, (1)--(5)]\label{thm:reflective-of-check}
\lean{ReflectZphi.Check.polyCheck, ReflectZphi.Check.admCheck, ReflectZphi.Check.winCheck, ReflectZphi.Check.chainCheck, ReflectZphi.Check.cellCheck, ReflectZphi.Check.reflective_of_check}\leanok
\uses{thm:bridge, lem:rootnorms, lem:cells, lem:split}
If the Boolean checks \texttt{polyCheck}, \texttt{admCheck}, \texttt{winCheck},
\texttt{chainCheck} and \texttt{cellCheck} all return \texttt{true} on the certificate of $G$,
and $G$ is over a PID with non-square discriminant, then the lattice is reflective, assuming
Vinberg's chamber theorem.
\end{theorem}
\begin{proof}\leanok
Theorem~\ref{thm:bridge}, with the one-sign certificate from Lemmas~\ref{lem:cells}
and~\ref{lem:split}.
\end{proof}

\begin{theorem}[the checker, per-norm form]\label{thm:reflective-of-checkM}
\lean{ReflectZphi.Sublat.reflective_of_checkM}\leanok
\uses{thm:bridge, lem:split}
The same conclusion from the per-norm sublattice form of the certificate.
\end{theorem}
\begin{proof}\leanok
As Theorem~\ref{thm:reflective-of-check}.
\end{proof}

\begin{theorem}[isometry certificates]\label{thm:iso}
\lean{ReflectZphi.checkIso, ReflectZphi.checkIso_sound}\leanok
If \texttt{checkIso}$(G_1,G_2,P)$ returns \texttt{true}, then $P\in GL_3(\Oo)$ and
$P^{\mathsf T}G_2P=G_1$.
\end{theorem}
\begin{proof}\leanok
The check is $P^{\mathsf T}G_2P=G_1$ and $N(\det P)=\pm1$.
\end{proof}

\chapter{The main theorem over \texorpdfstring{$\ZZ[\varphi]$}{Z[phi]}}\label{chap:main}

\begin{theorem}[paper, Thm.~12.1, reflectivity part]\label{thm:all-reflective}
\lean{ReflectZphi.Results.all_reflective, ReflectZphi.Results.classGram}\leanok
\uses{thm:reflective-of-check, thm:reflective-of-checkM, def:vinberg, thm:zphi}
Assume Vinberg's chamber theorem. For each $n=1,\dots,144$, the lattice $(\Oo^3,G_n)$ over
$\Oo=\ZZ[\varphi]$ is reflective, where $G_n$ is the Gram matrix of the representative of
class $\#n$ of the paper's Appendix~A.3.
\end{theorem}
\begin{proof}\leanok
For each class, the module \texttt{ReflectZphi.Data.Cert5.C\#\#\#} proves
\texttt{reflective (hV : VinbergChamber) : Reflective hd.le G} by
Theorem~\ref{thm:reflective-of-check} or~\ref{thm:reflective-of-checkM}, with every check
evaluated by \texttt{decide +kernel}. The file \texttt{ReflectZphi/Results.lean} collects
the $144$ theorems and fails to compile unless every theorem of the certificate library uses
only the axioms \texttt{propext}, \texttt{Classical.choice} and \texttt{Quot.sound}.
\end{proof}

\begin{theorem}[the isometries within the classes]\label{thm:iso-classes}
\uses{thm:iso}\leanok
Each of the $1140$ reflective candidate lattices over $\ZZ[\varphi]$ is isometric to the
representative of its class: the $996$ theorems \texttt{ReflectZphi.Data.Iso5\_\#\#\#.iso5\_\#\#\#\#\#}
are \texttt{checkIso} certificates, one for each member other than the representative.
Where Appendix~A.3 prints a member other than the representative, the comment of its entry
in \texttt{ReflectZphi/Results.lean} names the certificate.
\end{theorem}
\begin{proof}\leanok
Theorem~\ref{thm:iso}, evaluated by \texttt{decide +kernel}.
\end{proof}

\chapter{Proved in the paper only}\label{chap:not-formalised}

\begin{theorem}[transfer of reflectivity along an isometry]\label{thm:transfer}
\notready
\uses{def:reflective, thm:iso}
If $P^{\mathsf T}G_2P=G_1$ with $P\in GL_3(\Oo)$, then $(\Oo^3,G_1)$ is reflective if and only
if $(\Oo^3,G_2)$ is. (Elementary, but not formalised; Theorem~\ref{thm:all-reflective} is
therefore stated for the class representatives.)
\end{theorem}

\begin{theorem}[residue fields and admissible angles; paper, Prop.~9.1(3)--(5)]\label{thm:residues}
\notready
\uses{thm:zphi}
The residue fields of $\Oo$ at $2$, $3$ and $\pi$, and $\gamma(m)\in\Oo$ exactly for
$m\in\{2,3,4,5,6,10\}$.
\end{theorem}

\begin{theorem}[non-reflectivity certificates; paper, Prop.~7.6]\label{thm:nonrefl}
\notready
\uses{def:reflective}
A translation certificate shows that a candidate lattice is not reflective.
\end{theorem}

\begin{theorem}[completeness of the classification; paper, Sections~5--6 and Thm.~12.1]\label{thm:classification}
\notready
\uses{def:allcock, thm:nonrefl, thm:all-reflective, thm:transfer}
Every strongly squarefree reflective lattice over $\ZZ[\varphi]$ is isometric to exactly one of
the $144$ classes: the configurations and candidate lattices are complete (using Allcock's
thin-part theorem and the local-simplicity filter), every other candidate is
non-reflective, and the $144$ classes are pairwise non-isometric.
\end{theorem}

\begin{thebibliography}{9}
\bibitem{Vinberg} \`E.~B.~Vinberg, \emph{On groups of unit elements of certain quadratic
forms}, Mat. Sb. (N.S.) \textbf{87(129)} (1972), 18--36; English transl.: Math. USSR-Sb.
\textbf{16} (1972), 17--35.
\end{thebibliography}
